\subsection{Control, No-Load, and Load-Removal Tests}
\label{sec:method-control}

The final baseline controller acquired secondary-bulk voltage at PA1 through
continuous ADC DMA and controlled switching intervals through the MCU burst-enable
pin PA10. Saved UART settings, where available, identified the active
configuration rather than source-code defaults.

Loaded tests checked delivered voltage and burst structure at a fixed output setting. Because the records begin after startup, synchronized measurements of initial conditions, precharge, output rise and peak, the burst-enable signal (PA10), inhibition, decay, and fault status are still needed to quantify startup and shutdown.

Load-removal captures covered output voltage and switching around disconnection.
A complete assessment requires both the immediate stored-energy response and
later no-load cycling, with thresholds, latch state, and recovery recorded to
distinguish holdoff from a latched fault.

\begin{table}[htbp]
\centering
\caption{Signals needed for control-response tests.}
\label{tab:control-tests}
\begin{tabularx}{\textwidth}{@{}p{0.27\textwidth}X@{}}
\toprule
Condition & Observations\\
\midrule
Startup & Initial state, precharge, output rise, maximum voltage, and fault status\\
Loaded operation & Output voltage/ripple, load current, and PA10 pattern\\
Load removal & Immediate voltage excursion, enable response, and final no-load behavior\\
No load & Complete/partial burst structure and voltage cycling\\
Protection event & Calibrated threshold, inhibition, latch/holdoff state, and recovery\\
\bottomrule
\end{tabularx}
\end{table}
